\section{Which Checkpoints Get Selected?}
\label{sec:selection}

\Cref{sec:related} argued that ranking by criterion differs from ranking by
recency. That is a claim about the selected sets, not about accuracy, and it is
directly checkable. We inspect the selected epoch indices for the CIFAR-100
runs, for which per-seed selection records are available: six architectures
$\times$ five seeds $=30$ completed runs, hence 150 selected checkpoints at
$k{=}5$. \Cref{tab:selection} summarizes where they sit.

\begin{table}[t]
  \centering \small
  \setlength{\tabcolsep}{4pt}\renewcommand{\arraystretch}{1.12}
  \caption{Position of the selected checkpoints on CIFAR-100, averaged over
    five seeds. $T$ is the number of epochs actually completed (early stopping
    usually ends a run before its budget), $t^{*}$ the baseline's best epoch.
    \emph{Span} is $\max\mathcal{S}_5-\min\mathcal{S}_5+1$, so a contiguous
    block would span exactly 5. The last column counts members of
  $\mathcal{S}_5$ inside the final five epochs of the run.}
  \label{tab:selection}
  \resizebox{\linewidth}{!}{%
\begin{tabular}{@{}lccccc@{}}
    \toprule
    Model & $T$ & $t^{*}$ & \makecell{Span of\\$\mathcal{S}_5$}
    & \makecell{Non-contig.\\(of 5 seeds)}
    & $|\mathcal{S}_5\cap\text{last-}5|$ \\
    \midrule
    VGG16          & 75.8 & 63.8 & 17.6 & 5 & 1.0 \\
    DenseNet121    & 37.2 & 25.2 & 10.0 & 5 & 0.2 \\
    EfficientNetB0 & 23.8 & 11.8 &  9.0 & 5 & 0.2 \\
    MobileNetV2    & 50.0 & 38.0 & 16.0 & 5 & 0.4 \\
    ResNet101      & 46.6 & 34.6 & 14.4 & 5 & 1.6 \\
    ViT-Base       & 14.2 &  6.2 &  5.0 & 0 & 0.0 \\
    \midrule
    All            & 41.3 & 30.0 & 12.0 & 25 & 0.6 \\
    \bottomrule
  \end{tabular}}
\end{table}

\paragraph{The selected set is rarely a window.}
A block of five consecutive epochs spans exactly 5; the observed mean span is
12.0, and only 5 of 30 runs produce a contiguous set. In the other 25 the
criterion skips epochs, with a mean largest internal gap of 5.0 epochs. The
contiguous cases are all ViT-Base, which early-stops after roughly 14 epochs
and so has little room for its five best checkpoints to be anything but
adjacent---the exception that identifies the mechanism rather than one that
contradicts it. The longer the run, the further selection departs from a
window: VGG16 trains 75.8 epochs and spreads its five checkpoints over 17.6.

\paragraph{Nor is it concentrated at the end.}
Only 17 of the 150 selected checkpoints (11.3\%) lie in the final five epochs,
and in 20 of 30 runs $\mathcal{S}_5$ is entirely disjoint from that window. The
earliest selected checkpoint precedes the end of training by 16.4 epochs on
average, and $t^{*}$ itself sits at about 66\% of the completed trajectory---a
direct consequence of early stopping with patience 8--12. A last-$5$ average on
these same runs would therefore aggregate a largely different set of weights.
This establishes that the two rules are structurally distinct; it does not by
itself establish that ours is the better of the two, which would require the
matched comparison we flag as future work below.

\paragraph{Storage.}
Runs retain 41.3 checkpoints on average under naive retention. As noted in
\cref{sec:method}, a running top-$k$ set reduces this to $k{+}1$ states with an
identical result, so the practical cost is six checkpoints at $k{=}5$.

\section{Discussion}
\label{sec:discussion}

\paragraph{What the matched-trajectory design buys.}
Because the baseline and every $k$ come from one run per seed, the comparison
controls for architecture, initialization, data order, optimizer, schedule and
budget, so the consistent direction of the effect across three tasks, eight models and 34
configurations is hard to attribute to training variance. A plausible mechanism
is that the single best checkpoint's advantage is partly transient---it won a
noisy comparison---and that averaging several states of comparable validation
quality keeps what they agree on while cancelling part of what is idiosyncratic
to each.

\paragraph{Limitations.}
Three are worth stating plainly. \emph{First}, we establish an advantage over
single-checkpoint selection, not over alternative aggregation strategies.
Last-$k$ averaging, EMA, post-hoc SWA and prediction ensembling are not run
here under a matched protocol; \cref{sec:selection} shows last-$k$ would select
different checkpoints but not that it would perform worse. \emph{Second}, the
averaged model receives a BN recalibration pass that the single-checkpoint
baseline does not need and does not get. For architectures with normalization
layers this leaves open how much of the gain is attributable to averaging
versus to recalibration; the clean control is a BN-recalibrated baseline, which
we regard as the most important missing ablation. \emph{Third}, the selection
analysis of \cref{sec:selection} rests on CIFAR-100 records only, so its
quantitative constants should not be read as universal, though the qualitative
conclusion follows from early stopping being active, which is common.
